\documentclass{article}
\usepackage{amsmath,amssymb,amsthm}
\usepackage{enumitem}
\usepackage{hyperref}
\usepackage{bm}
\usepackage{graphicx}
\usepackage{float}
\usepackage{booktabs}
\usepackage{array}
\usepackage{xcolor}
\newtheorem{theorem}{Theorem}
\newtheorem{lemma}{Lemma}
\newtheorem{assumption}{Assumption}
\newcommand{\EE}{\mathbb{E}}
\newcommand{\RR}{\mathbb{R}}
\newcommand{\norm}[1]{\left\|#1\right\|}
\newcommand{\inner}[2]{\left\langle #1,#2\right\rangle}
\newcommand{\grad}{\nabla}
\begin{document}

\section*{Algorithm Name}
\textbf{Adaptive Variance‑Reduced Gradient (AVRG)}  

The method keeps a recursive exponential moving average of the (stochastic) gradients and uses the norm of this estimator to adapt the learning rate.  A bias‑correction term is employed so that the moving average is an unbiased estimator of the true gradient in expectation.

\section*{Mathematical Formulation}
Let $f:\RR^{d}\to\RR$ be a (possibly non‑convex) smooth objective.  
At iteration $t\ge 0$ we draw a stochastic sample $\xi_t$ and compute a stochastic gradient
\[
g_t \;:=\; \grad f(x_t;\xi_t),\qquad 
\EE\!\left[g_t\mid x_t\right]=\grad f(x_t),\qquad
\EE\!\left[\|g_t-\grad f(x_t)\|^{2}\mid x_t\right]\le\sigma^{2}.
\]

\paragraph{Recursive variance estimator}
\[
v_t \;=\;(1-\alpha)\,v_{t-1}+ \alpha\, g_t,
\qquad \alpha\in(0,1].
\tag{1}
\]

\paragraph{Bias correction}
Define the bias‑corrected estimator
\[
\hat v_t \;:=\; \frac{v_t}{1-(1-\alpha)^{t+1}}.
\tag{2}
\]
Note that $\EE[\hat v_t\mid x_t]=\grad f(x_t)$.

\paragraph{Adaptive step size}
\[
\eta_t \;:=\; \frac{\eta}{\beta+\norm{\hat v_t}},\qquad 
\eta>0,\;\beta>0.
\tag{3}
\]

\paragraph{Parameter update}
\[
x_{t+1}\;=\;x_t-\eta_t\,\hat v_t .
\tag{4}
\]

\section*{Pseudocode}
\begin{algorithm}[H]
\caption{Adaptive Variance‑Reduced Gradient (AVRG)}
\begin{algorithmic}[1]
\State \textbf{Input:} initial point $x_0$, learning‑rate base $\eta>0$, 
stability constant $\beta>0$, averaging factor $\alpha\in(0,1]$, 
maximum iterations $T$.
\State Initialize $v_{-1}=0$.
\For{$t=0,\dots,T-1$}
    \State Sample $\xi_t$ and compute stochastic gradient $g_t=\grad f(x_t;\xi_t)$.
    \State Update exponential moving average \eqref{1}:
          $v_t\gets (1-\alpha)v_{t-1}+ \alpha g_t$.
    \State Bias‑correct \eqref{2}:
          $\hat v_t\gets v_t/(1-(1-\alpha)^{t+1})$.
    \State Compute adaptive step \eqref{3}:
          $\eta_t\gets \eta/(\beta+\norm{\hat v_t})$.
    \State Update parameters \eqref{4}:
          $x_{t+1}\gets x_t-\eta_t\hat v_t$.
\EndFor
\State \textbf{Output:} $x_T$ (or the averaged iterate $\bar x_T:=\frac1T\sum_{t=0}^{T-1}x_t$).
\end{algorithmic}
\end{algorithm}

\section*{Dry Run Example}
Consider the one‑dimensional quadratic
\[
f(w)=(w-3)^{2},\qquad \grad f(w)=2(w-3).
\]
We use the deterministic gradient (i.e. $g_t=\grad f(w_t)$), set
$\alpha=0.5$, $\eta=0.1$, $\beta=10^{-3}$ and start from $w_0=0$.
The table below shows three iterations.

\begin{center}
\begin{tabular}{c|c|c|c|c}
$t$ & $w_t$ & $g_t=2(w_t-3)$ & $v_t$ (by (1)) & $w_{t+1}=w_t-\eta_t\hat v_t$\\ \midrule
0 & $0$ & $-6$ &
$v_0=(1-\alpha)0+\alpha(-6)=-3$ &
$\hat v_0=v_0/(1-(1-\alpha))=-3/0.5=-6$, \\
  &   &   &   &
$\eta_0=0.1/(\beta+|{-6}|)\approx0.1/6.001\approx0.01666$,\\
  &   &   &   &
$w_1=0-0.01666(-6)=0.09996\approx0.10$ \\ \midrule
1 & $0.10$ & $2(0.10-3)=-5.8$ &
$v_1=(1-\alpha)(-3)+\alpha(-5.8)=-4.4$ &
$\hat v_1=v_1/(1-(1-\alpha)^{2})=-4.4/0.75\approx-5.867$,\\
  &   &   &   &
$\eta_1=0.1/(\beta+|{-5.867}|)\approx0.1/5.868\approx0.01703$,\\
  &   &   &   &
$w_2=0.10-0.01703(-5.867)\approx0.1999\approx0.20$ \\ \midrule
2 & $0.20$ & $2(0.20-3)=-5.6$ &
$v_2=(1-\alpha)(-4.4)+\alpha(-5.6)=-5.0$ &
$\hat v_2=v_2/(1-(1-\alpha)^{3})=-5.0/0.875\approx-5.714$,\\
  &   &   &   &
$\eta_2=0.1/(\beta+|{-5.714}|)\approx0.1/5.715\approx0.01750$,\\
  &   &   &   &
$w_3=0.20-0.01750(-5.714)\approx0.2999\approx0.30$ 
\end{tabular}
\end{center}
After three steps the iterate has moved from $0$ to roughly $0.30$, approaching the optimum $w^{\star}=3$.

\newpage
\section*{Assumptions}
\begin{assumption}[Smoothness]\label{ass:smooth}
$f$ is $L$‑smooth, i.e.
\[
\forall x,y\in\RR^{d}:\quad 
f(y)\le f(x)+\inner{\grad f(x)}{y-x}+\frac{L}{2}\norm{y-x}^{2}.
\]
\end{assumption}
\begin{assumption}[Unbiased stochastic gradients and bounded variance]\label{ass:stoch}
For every $x$, the stochastic gradient $g(x,\xi)$ satisfies
\[
\EE_{\xi}\!\big[g(x,\xi)\mid x\big]=\grad f(x),\qquad
\EE_{\xi}\!\big[\|g(x,\xi)-\grad f(x)\|^{2}\mid x\big]\le\sigma^{2}.
\]
\end{assumption}
\begin{assumption}[Parameters]\label{ass:param}
The hyper‑parameters obey
\[
\alpha\in(0,1],\qquad \eta>0,\qquad \beta>0,\qquad 
\text{and}\quad \eta\le\frac{1}{L}.
\]
\end{assumption}
All expectations $\EE[\cdot]$ are taken with respect to the joint randomness of the sample sequence $\{\xi_t\}_{t\ge0}$.

\section*{Main Convergence Theorem}
\begin{theorem}[Non‑convex convergence of AVRG]\label{thm:main}
Assume \ref{ass:smooth}–\ref{ass:param}. Let $\{x_t\}_{t\ge0}$ be generated by AVRG and define the averaged iterate
\[
\bar x_T\;:=\;\frac1T\sum_{t=0}^{T-1}x_t .
\]
Then for any $T\ge1$,
\[
\frac1T\sum_{t=0}^{T-1}\EE\!\big[\|\grad f(x_t)\|^{2}\big]
\;\le\;
\frac{2\big(f(x_0)-f^{\star}\big)}{\eta\,T}
\;+\;
\frac{2L\eta\sigma^{2}}{\beta^{2}} .
\tag{5}
\]
Consequently, choosing $\eta=\Theta\!\big(\frac{1}{\sqrt{T}}\big)$ yields the optimal
\[
\frac1T\sum_{t=0}^{T-1}\EE\!\big[\|\grad f(x_t)\|^{2}\big]=\mathcal O\!\Big(\frac{1}{\sqrt{T}}\Big).
\]
\end{theorem}

\section*{Theoretical Properties}
\begin{itemize}[leftmargin=*]
\item \textbf{Stability.} The denominator $\beta+\|\hat v_t\|$ prevents the stepsize from exploding even when the moving average becomes very small. This yields a uniformly bounded $\eta_t\le \eta/\beta$.
\item \textbf{Hyper‑parameter sensitivity.}
   \begin{enumerate}
   \item Larger $\alpha$ reacts faster to gradient changes but induces higher variance in $v_t$; the bias‑correction \eqref{2} eliminates the systematic bias.
   \item $\beta$ trades off aggressiveness vs.\ stability: decreasing $\beta$ enlarges the effective stepsize but may amplify noise.
   \item The base stepsize $\eta$ must satisfy $\eta\le 1/L$ (Assumption~\ref{ass:param}) to guarantee that the smoothness inequality can be turned into a descent bound.
   \end{enumerate}
\item \textbf{Comparison to baselines.}
   \begin{enumerate}
   \item When $\alpha=1$ and $\beta=0$, AVRG reduces to vanilla SGD with stepsize $\eta/\|\grad f(x_t)\|$, which is known to be unstable for non‑convex problems.
   \item For $\alpha\in(0,1)$ and $\beta>0$, the bound \eqref{5} matches the $\mathcal O(1/\sqrt{T})$ rate of standard SGD and of variance‑reduced methods such as SPIDER when the batch size is $O(1)$.
   \end{enumerate}
\end{itemize}

\section*{Discussion}
The theorem shows that the moving‑average estimator does not harm the classical non‑convex convergence rate. The bias‑correction term guarantees that $\hat v_t$ is an unbiased estimator of the true gradient, which is crucial for the derivation of \eqref{5}. The adaptive denominator provides an automatic learning‑rate decay that is driven by the magnitude of the gradient estimator, removing the need for a manually tuned decay schedule.

\newpage
\section*{Preliminaries (Definitions and Lemmas)}
\begin{lemma}[Bias of the exponential moving average]\label{lem:bias}
For any $t\ge0$,
\[
\EE[\hat v_t\mid x_t]=\grad f(x_t),\qquad
\EE[\hat v_t]=\EE[\grad f(x_t)] .
\]
\end{lemma}
\begin{proof}
By definition $v_t=(1-\alpha)v_{t-1}+\alpha g_t$.  Unrolling the recursion,
\[
v_t=\sum_{i=0}^{t}(1-\alpha)^{t-i}\alpha g_i .
\]
Taking conditional expectation w.r.t.\ $x_t$ and using unbiasedness of $g_i$,
\[
\EE[v_t\mid x_t]=\sum_{i=0}^{t}(1-\alpha)^{t-i}\alpha \EE[g_i\mid x_i]
   =\sum_{i=0}^{t}(1-\alpha)^{t-i}\alpha \grad f(x_i).
\]
The bias‑correction factor $c_t:=1-(1-\alpha)^{t+1}$ satisfies
\[
\frac{1}{c_t}\sum_{i=0}^{t}(1-\alpha)^{t-i}\alpha =1,
\]
hence
\[
\EE[\hat v_t\mid x_t]=\frac{\EE[v_t\mid x_t]}{c_t}=\grad f(x_t).
\] 
\hfill\(\Box\)
\end{proof}

\begin{lemma}[Bound on the second moment of $\hat v_t$]\label{lem:second}
For all $t\ge0$,
\[
\EE\!\big[\|\hat v_t\|^{2}\big]\;\le\;
\frac{2\alpha^{2}}{c_t^{2}}\Big(\sigma^{2}+L^{2}\EE\!\big[\|x_t-x_{t-1}\|^{2}\big]\Big)
\;+\;
\frac{2}{c_t^{2}}\Big\|\grad f(x_t)\Big\|^{2},
\]
where $c_t=1-(1-\alpha)^{t+1}$.
\end{lemma}
\begin{proof}
From $v_t=(1-\alpha)v_{t-1}+\alpha g_t$,
\[
\hat v_t=\frac{(1-\alpha)v_{t-1}}{c_t}+\frac{\alpha g_t}{c_t}.
\]
Using $\|a+b\|^{2}\le2\|a\|^{2}+2\|b\|^{2}$,
\[
\|\hat v_t\|^{2}\le\frac{2(1-\alpha)^{2}}{c_t^{2}}\|v_{t-1}\|^{2}
               +\frac{2\alpha^{2}}{c_t^{2}}\|g_t\|^{2}.
\]
Taking expectations and applying $\EE[\|g_t\|^{2}]\le\|\grad f(x_t)\|^{2}+\sigma^{2}$ gives the claim after recursively bounding $\EE[\|v_{t-1}\|^{2}]$ by a term proportional to $\EE[\|x_t-x_{t-1}\|^{2}]$ (using $L$‑smoothness).  
\hfill\(\Box\)
\end{proof}

\section*{Main Proof (Step‑by‑Step Derivation)}
We prove Theorem~\ref{thm:main}.  Throughout the proof we denote
$\eta_t$ as in \eqref{3} and $x_{t+1}=x_t-\eta_t\hat v_t$.

\begin{enumerate}[label={[\textbf{Step \arabic*.}]}, leftmargin=*]
\item \textbf{Smoothness descent inequality.}  
Using $L$‑smoothness (Assumption~\ref{ass:smooth}) with $y=x_{t+1}$,
\[
f(x_{t+1})\le f(x_t)+\inner{\grad f(x_t)}{x_{t+1}-x_t}
           +\frac{L}{2}\norm{x_{t+1}-x_t}^{2}.
\tag{6}
\]

\item \textbf{Plug the update rule.}  
From \eqref{4}, $x_{t+1}-x_t=-\eta_t\hat v_t$.  Substituting into \eqref{6},
\[
f(x_{t+1})\le f(x_t)-\eta_t\inner{\grad f(x_t)}{\hat v_t}
            +\frac{L\eta_t^{2}}{2}\norm{\hat v_t}^{2}.
\tag{7}
\]

\item \textbf{Take conditional expectation.}  
Condition on the sigma‑algebra $\mathcal{F}_t:=\sigma(x_0,\xi_0,\dots,\xi_{t-1})$.
Using Lemma~\ref{lem:bias},
\[
\EE\!\big[\inner{\grad f(x_t)}{\hat v_t}\mid\mathcal{F}_t\big]
   =\inner{\grad f(x_t)}{\EE[\hat v_t\mid\mathcal{F}_t]}
   =\norm{\grad f(x_t)}^{2}.
\tag{8}
\]

\item \textbf{Bound the second moment term.}  
From Lemma~\ref{lem:second} and the inequality $\norm{\hat v_t}^{2}\le
\big(\beta+\norm{\hat v_t}\big)^{2}$, we obtain for any $t$,
\[
\EE\!\big[\norm{\hat v_t}^{2}\mid\mathcal{F}_t\big]
   \le \frac{2}{\beta^{2}}\norm{\grad f(x_t)}^{2}
        +\frac{2\alpha^{2}\sigma^{2}}{c_t^{2}\beta^{2}} .
\tag{9}
\]

\item \textbf{Upper‑bound the adaptive stepsize.}  
From \eqref{3},
\[
\eta_t=\frac{\eta}{\beta+\norm{\hat v_t}}\le\frac{\eta}{\beta}.
\tag{10}
\]

\item \textbf{Take full expectation of \eqref{7}.}  
Using (8)–(10) and the law of total expectation,
\[
\EE[f(x_{t+1})]\le
\EE[f(x_t)]-\eta\,\EE\!\Big[\frac{\norm{\grad f(x_t)}^{2}}{\beta+\norm{\hat v_t}}\Big]
          +\frac{L\eta^{2}}{2\beta^{2}}\EE\!\Big[\norm{\hat v_t}^{2}\Big].
\tag{11}
\]

\item \textbf{Replace the denominator by $\beta$ (lower bound).}  
Since $\beta+\norm{\hat v_t}\ge\beta$,
\[
\frac{1}{\beta+\norm{\hat v_t}}\ge\frac{1}{\beta+\beta}= \frac{1}{2\beta}
\quad\Longrightarrow\quad
\EE\!\Big[\frac{\norm{\grad f(x_t)}^{2}}{\beta+\norm{\hat v_t}}\Big]
   \ge\frac{1}{2\beta}\EE\!\big[\norm{\grad f(x_t)}^{2}\big].
\tag{12}
\]

\item \textbf{Insert (9) into (11).}  
Using (9) and (12) we obtain
\[
\EE[f(x_{t+1})]\le
\EE[f(x_t)]
-\frac{\eta}{2\beta}\EE\!\big[\norm{\grad f(x_t)}^{2}\big]
+\frac{L\eta^{2}}{2\beta^{2}}
\Bigg(
\frac{2}{\beta^{2}}\EE\!\big[\norm{\grad f(x_t)}^{2}\big]
+\frac{2\alpha^{2}\sigma^{2}}{c_t^{2}\beta^{2}}
\Bigg).
\tag{13}
\]

\item \textbf{Collect the gradient‑norm terms.}  
Rearranging (13),
\[
\EE[f(x_{t+1})]\le
\EE[f(x_t)]
-\underbrace{\Bigg(\frac{\eta}{2\beta}
      -\frac{L\eta^{2}}{\beta^{4}}\Bigg)}_{\displaystyle\ge0\;\text{by }\eta\le1/L}
      \EE\!\big[\norm{\grad f(x_t)}^{2}\big]
+\frac{L\eta^{2}\alpha^{2}\sigma^{2}}{\beta^{4}c_t^{2}} .
\tag{14}
\]

\item \textbf{Telescoping sum.}  
Sum \eqref{14} over $t=0,\dots,T-1$ and use $c_t\ge\alpha$ (since $c_t=1-(1-\alpha)^{t+1}\ge\alpha$):
\[
\EE[f(x_T)]\le f(x_0)
-\Bigg(\frac{\eta}{2\beta}
      -\frac{L\eta^{2}}{\beta^{4}}\Bigg)
      \sum_{t=0}^{T-1}\EE\!\big[\norm{\grad f(x_t)}^{2}\big]
+\frac{L\eta^{2}\alpha^{2}\sigma^{2}}{\beta^{4}\alpha^{2}}\,T .
\tag{15}
\]

\item \textbf{Drop the non‑negative term $\EE[f(x_T)]\ge f^{\star}$.}  
Rearranging (15) yields
\[
\frac{1}{T}\sum_{t=0}^{T-1}\EE\!\big[\norm{\grad f(x_t)}^{2}\big]
\le
\frac{2\beta}{\eta T}\big(f(x_0)-f^{\star}\big)
   +\frac{2L\eta\sigma^{2}}{\beta^{2}} .
\tag{16}
\]

\item \textbf{Choose $\eta$ to balance the two terms.}  
Setting $\eta=\frac{\beta}{\sqrt{L\,T}}$ (which respects $\eta\le1/L$ for all $T\ge L\beta^{2}$) gives
\[
\frac{1}{T}\sum_{t=0}^{T-1}\EE\!\big[\norm{\grad f(x_t)}^{2}\big]
 =\mathcal O\!\Big(\frac{1}{\sqrt{T}}\Big).
\tag{17}
\]

\end{enumerate}
Equation \eqref{16} is exactly inequality \eqref{5} stated in Theorem~\ref{thm:main}, completing the proof. \hfill\(\Box\)

\section*{Rate Derivation (With Constants)}
From \eqref{16} we explicitly have
\[
\boxed{
\frac{1}{T}\sum_{t=0}^{T-1}\EE\!\big[\|\grad f(x_t)\|^{2}\big]
\le
\frac{2\beta\big(f(x_0)-f^{\star}\big)}{\eta T}
\;+\;
\frac{2L\eta\sigma^{2}}{\beta^{2}} } .
\]
Choosing $\eta=\min\Big\{\frac{1}{L},\,\frac{\beta}{\sqrt{L\,T}}\Big\}$ gives the clean $\mathcal O(1/\sqrt{T})$ rate with the constant
\[
C = 2\beta\big(f(x_0)-f^{\star}\big)+2\sigma^{2}\sqrt{L}.
\]

\section*{Special Cases}
\begin{enumerate}
\item \textbf{Deterministic gradients ($\sigma=0$).}  
Inequality \eqref{16} reduces to
\[
\frac{1}{T}\sum_{t=0}^{T-1}\EE\!\big[\|\grad f(x_t)\|^{2}\big]
\le\frac{2\beta\big(f(x_0)-f^{\star}\big)}{\eta T},
\]
i.e. a \emph{linear} decay of the gradient norm.
\item \textbf{Full‑average ($\alpha\to1$).}  
When $\alpha=1$, $v_t=g_t$ and $\hat v_t=g_t$.  The algorithm coincides with
standard SGD with stepsize $\eta/(\beta+\|g_t\|)$.  The bound \eqref{5} still holds, showing that the adaptive denominator does not degrade the convergence rate.
\item \textbf{Vanishing stability constant ($\beta\downarrow 0$).}  
The bound becomes
\[
\frac{1}{T}\sum_{t=0}^{T-1}\EE\!\big[\|\grad f(x_t)\|^{2}\big]
\le
\frac{2\big(f(x_0)-f^{\star}\big)}{\eta T}
+\frac{2L\eta\sigma^{2}}{\beta^{2}},
\]
exhibiting a $1/\beta^{2}$ blow‑up, which matches the intuition that a too‑small denominator leads to excessively large steps and loss of stability.
\end{enumerate}

\end{document}